\documentclass[10pt,a4paper]{amsart}
\usepackage[margin=19mm]{geometry}
\usepackage{amsmath,amssymb,mathtools}
\usepackage[scr=rsfs,cal=euler]{mathalfa}
\usepackage{xfrac}
\usepackage[unicode,hidelinks,bookmarks=false]{hyperref}
\newcommand{\ie}{i.e.\ }
\newcommand{\cf}{cf.\ }
\newcommand{\resp}{respectively\ }
\newcommand{\Subsection}{\subsection}
\providecommand{\nobreakdash}{\nobreak\hbox{-}}
\setlength{\emergencystretch}{2em}
\multlinegap0pt
\usepackage{bm}
\usepackage{bbm}
\usepackage{dsfont}
\usepackage{xspace}
\usepackage{cancel}
\usepackage{mathtools}
\usepackage{amsthm}
\makeatletter
\@ifundefined{theorem}{\newtheorem{theorem}{Theorem}}{}
\@ifundefined{proposition}{\newtheorem{proposition}[theorem]{Proposition}}{}
\makeatother
\renewcommand{\bar}{\overline}
\title[Bidihedral skew braces]{Bidihedral skew braces}
\author{Alan Koch, Paul J. Truman}
\date{Source: arXiv:2607.15917v1 (CC BY 4.0)}
\begin{document}
\maketitle

\noindent\emph{Source and licence.} Bidihedral skew braces. Version arXiv:2607.15917v1 (CC BY 4.0), distributed under the Creative Commons Attribution 4.0 International licence, as recorded in the arXiv licence metadata for this exact version and shown on the abstract page https://arxiv.org/abs/2607.15917v1. Full rights notice: licenses/arxiv-2607.15917-CC-BY-4.0.txt.\par\smallskip

\noindent\textit{Editorial scope.} Counted unit: the Proposition whose source label is \texttt{prop\_di\_cyc\_realise} in the article (source0.tex lines 571--575, source label \texttt{prop\_di\_cyc\_realise}), delivered together with the article's own complete original proof of it (source0.tex lines 576--618, one \texttt{proof} environment of 43 lines). No mathematics is written, repaired or re-derived: statement and proof are the article's own text line for line. Required context retained uncounted: prop\_bicylic\_necessary (lines 182--199); prop\_d\_n\_realise (lines 233--237); prop\_d\_n\_isomorphism (lines 276--278); prop\_d\_n\_2\_realise (lines 333--337); prop\_di\_cyc\_necessary (lines 525--534). Every definition, display and label of those lines is retained unchanged. Omitted from this package and not redistributed: 1--181, 200--232, 238--275, 279--332, 338--524, 535--570, 619--763. Nothing inside a retained excerpt is omitted or altered: every retained body is delivered exactly as the article's lines are written, so no omission receipt of any kind accompanies this package. Citations: the retained excerpts carry no citation command at all, so no bibliography entry of the article is transcribed and no reference is redistributed. Reference adaptation: none was needed and none was made; every reference inside the delivered unit resolves inside it, so no unresolved reference or citation is produced. Printed slips kept literal: no printed word, symbol, hypothesis or number of the counted statement or of its proof was altered. Layout and class: the article's own class is \texttt{amsart} with packages including amssymb, amsthm, amsmath, xypic, mathrsfs, lineno, color, ifthen, xcolor; the wrapper is the lane's pinned \texttt{amsart} editorial wrapper at 10pt on A4, which restates the theorem-like environments the retained text needs and adds the article's own macro definitions that the retained text needs (\texttt{\textbackslash bar}), with extra wrapper packages (bm, bbm, dsfont, xspace, cancel, mathtools). The delivered numbering is the unit's own. Assets: no retained span contains a figure environment, a graphics-inclusion command, a PDF-page inclusion, a table or a TikZ picture; no image file is copied into this package, the package declares no asset dependency and no image file is redistributed with it. Licence: version arXiv:2607.15917v1 of this article is distributed under the Creative Commons Attribution 4.0 International licence, as recorded in the arXiv licence metadata for this exact version and shown on the abstract page https://arxiv.org/abs/2607.15917v1. A full rights notice is delivered at licenses/arxiv-2607.15917-CC-BY-4.0.txt. Characters: every retained excerpt is pure ASCII, so the delivered engine, XeLaTeX with its default Latin Modern text font, reports no missing character.
\begin{proposition} \label{prop_bicylic_necessary} 
With the notation above, we have the following:
\begin{enumerate}
\item[]
\item The integer $ d $ satisfies
\[ d= \begin{cases} 
n \mbox{ or } n/2 & \mbox{if } 8 \mid n \\
n & \mbox{otherwise;} \end{cases} \]
\item The $ \gamma $-function of $ G $ satisfies
\begin{align*}
\gamma_{s}(r) &= r^{k} \mbox{ with } k^{2} \equiv 1 \pmod{n} \\
\gamma_{s}(s) &= s \\
\gamma_{r}(r) &= r^{1+d} \\
\gamma_{r}(s) &= r^{(1+d)k-1}s,
\end{align*}
and is determined by these values. 
\end{enumerate}
\end{proposition}

\begin{proposition} \label{prop_d_n_realise}
Let $ k \in K_{n} $ and define $ \circ $ on $ G^{\bullet} $ by
\[ r^{i}s^{u} \circ r^{j}s^{v} = \left( r^{k^{v}i}s^{u} \right)\left( r^{k^{u}j}s^{v} \right). \]
Then $ \mathscr{A}_{k} = (G,\cdot,\circ) $ is a skew brace that realises the conditions described in Proposition \ref{prop_bicylic_necessary} with $ d=n $. 
\end{proposition}

\begin{proposition} \label{prop_d_n_isomorphism}
For $ k \in K_{n} $, every automorphism of $ G^{\bullet} $ is a skew brace automorphism of $ \mathscr{A}_{k} $. 
\end{proposition}

\begin{proposition} \label{prop_d_n_2_realise}
Suppose that $ 8 \mid n $. Let $ k \in K_{n} $ and let $ \circ $ be as described in Proposition \ref{prop_d_n_realise}. Let $ c=r^{n/2} $, and define $ \star $ on $ G $ by
\[ r^{i}s^{u} \star r^{j}s^{v} = \left( r^{i}s^{u} \circ r^{j}s^{v} \right) \circ c^{i(j+v)}. \]
Then $ \mathscr{B}_{k} = (G,\cdot,\star) $ is a skew brace that realises the conditions described in Proposition \ref{prop_bicylic_necessary} with $ d=n/2 $. 
\end{proposition}

\begin{proposition} \label{prop_di_cyc_necessary}
The $ \gamma $-function satisfies
\begin{align*}
\gamma_{a}(r) &= r^{k} \mbox{ with } k^{2} \equiv 1 \pmod{n} \\
\gamma_{a}(a) &= r^{-2}a \\
\gamma_{r}(r) &= r^{-1} \\
\gamma_{r}(a) &= r^{k-1}a,
\end{align*}
and is determined by these values. 
\end{proposition}

\begin{proposition} \label{prop_di_cyc_realise}
Let $ k \in K_{n} $ and consider the skew brace $ \mathscr{A}_{-k} = (G,\cdot,\circ) $. Define $ \ast  $ on $ G^{\bullet} $ by
\[ r^{i}s^{u} \ast r^{j}s^{v} = \left( r^{i}s^{u} \right) \circ \left( r^{(-1)^{i}j}s^{v} \right). \]
Then $ \mathscr{C}_{k} = (G,\cdot,\ast) $ is a skew brace that realises the conditions described in Proposition \ref{prop_di_cyc_necessary} with $ a=r^{-k}s $. 
\end{proposition}

\begin{proof}
By Proposition \ref{prop_d_n_isomorphism} the automorphism $ \beta_{-1} $ of $ G^{\bullet} $ is also an automorphism of $ G^{\circ} $; using this automorphism we may rewrite the definition of $ \ast $ as 
\[ r^{i}s^{u} \ast r^{j}s^{v} = \left( r^{i}s^{u} \right) \circ \beta_{-1}^{i}\left( r^{j}s^{v} \right). \]
We see immediately that $ e $ is an identity element with respect to $ \ast $ and that inverses are given by $ \bar{r^{i}} = r^{(-1)^{i+1}i} $ and $ \bar{r^{i}s} = r^{(-1)^{i}}s $. 

As in the proof of Proposition \ref{prop_d_n_2_realise}, we show $ \ast $ is associative on $ G $, that $ G^{\ast} $ is dihedral, and that the brace relation is satisfied by exploiting the analogous properties for $ G^{\circ} $. Let $ x=r^{i}s^{u}, y=r^{j}s^{v}, $ and $ x = r^{\ell}s^{w} $. Then
\begin{align*}
x \ast (y \ast z) &= x \ast (y \circ \beta_{-1}^{j}(z)) \\
&= x \circ \beta_{-1}^{i}(y \circ \beta_{-1}^{j}(z)) \\
&= x \circ \beta_{-1}^{i}(y) \circ \beta_{-1}^{i+j}(z), 
\end{align*}
and
\begin{align*}
(x \ast y) \ast z &= (x \circ \beta_{-1}^{i}(y)) \ast z \\
&= x \circ \beta_{-1}^{i}(y) \circ \beta_{-1}^{i+(-1)^{i}j}(z). 
\end{align*}
Since $ \beta_{-1} $ has order $ 2 $ we have $ \beta_{-1}^{i+j} = \beta_{-1}^{i+(-1)^{i}j} $; it follows that $ \ast $ is associative on $ G $, and so $ G^{\ast} $ is a group. To show that it is dihedral, we note first that 
\[ rs \ast rs = rs \circ r^{-1}s = r^{-k}sr^{k}s = r^{-2k}, \]
and that (since $ (k,n)=1 $) $ r^{-2k} $ has order $ n/2 $ in $ G^{\ast} $. Therefore $ rs $ has order $ n $ in $ G^{\ast} $. It is immediate from the definition of $ \ast $ that $ r $ has order $ 2 $ in $ G^{\ast} $. Now we have
\[ r \ast rs = r \circ r^{-1}s = r^{-k}r^{-1}s = r^{-1-k}s, \]
and
\[ \bar{rs} \ast r = r^{-1}s \ast r = r^{-1}s \circ r^{-1} = r^{-1}sr^{k} = r^{-1-k}s. \]
Hence $ r \ast rs = \bar{rs} \ast r $, and so $ G^{\ast} $ is indeed dihedral. 

Finally, to show that the brace relation is satisfied we exploit the fact that $ \mathscr{A}_{-k} $ is a skew brace. We have
\begin{align*}
x \ast (yz) &= x \circ \beta_{-1}^{i}(yz) \\
&= x \circ \beta_{-1}^{i}(y)\beta_{-1}^{i}(z) \\
&= (x \circ \beta_{-1}^{i}(y))x^{-1}(x \circ \beta_{-1}^{i}(z)) \\
&= (x \ast y)x^{-1}(x \ast z). 
\end{align*}
Hence $ (G,\cdot,\ast) $ is a skew brace. We verify that it realises the conditions described in Proposition \ref{prop_di_cyc_necessary} with $ a=r^{-k}s $. Recall that we have set $ (G,\cdot,\circ) = \mathscr{A}_{-k} $.
\begin{align*}
\gamma_{a}(r) &= (r^{-k}s)^{-1}(r^{-k}s \ast r) = (r^{-k}s)(r^{-k}s \circ r^{-1}) = r^{-k}s r^{-k} s r^{-(-k)} = r^{k} \\
%
\gamma_{a}(a) &= (r^{-k}s)^{-1}(r^{-k}s \ast r^{-k}s) = (r^{-k}s)(r^{-k}s \circ r^{k}s)= (r^{-k}s)(r^{(-k)^{2}}s r^{-k^{2}}s) \\
& = r^{-k-2}s = r^{-2} a \\
%
\gamma_{r}(r) &= r^{-1}(r \ast r) = r^{-1}(r \circ r^{-1}) = r^{-1} \\
%
\gamma_{r}(a) &= r^{-1}(r \ast r^{-k}s) = r^{-1}( r^{-k} r^{k}s ) = r^{-1}s = r^{k-1}a.
\end{align*}
\end{proof}

\end{document}
